\scanpage{36}
L'élément $\alpha\in P(\mathbb C)$ est relié à $f_\infty$ par
\[
\overline\alpha=f_\infty\alpha
\]
donc \deleted{$f_\infty$} à priori
\[
f_\infty=1\Longleftrightarrow\alpha\in P(\mathbb R)\quad\text{i.e.}\quad\mathbb R(\alpha)=\mathbb R
\]

\deleted{Comme on a la relation $f_\infty j_1 f_\infty^{-1}=j_2$, cette condition
implique $j_1=j_2$ donc $j_1=j_2=\frac12 i$, i.e. $\unsure{G_2|G}$ [si on
admet conj. de Hodge] $G^0\xrightarrow[\sim]{\varepsilon}G_m$}

Or on vérifie aussitôt
\[
\boxed{\varepsilon(f_\infty)=-1}
\]
ce qui nous montre qu'on a tjrs
\[
\boxed{f_\infty\ne1\quad\text{i.e.}\quad\mathbb R(\alpha)=\mathbb C}
\quad\text{(si $M$ contient le motif de Tate)}
\]

Si $M$ ne contient pas le motif de Tate, on voit
que $f_\infty=1\Longrightarrow j_1=j_2$ i.e. $j_1=j_2=\frac12 i$, donc en
vertu de la conj. de Hodge, il faut que
$G^0=\mathrm{Im}\,j_1$, donc ou bien $G^0=\{e\}$ \deleted{(\unsure{motifs de}
\unsure{dimension} $0$, alors $G=\{e\}$, \inserted{ici $G=G^0$, car} $G=\mathbb Z/2\mathbb Z$, dans ce
dernier cas $f_\infty$} \inserted{donc $G=1$}
[\uline{NB} (si $G$ est disconnexe, on sait que
$f_\infty\notin G^0(\mathbb Q)$)], ou bien $G^{\deleted{\unsure{0}}}=G_m/\mu_n$ \inserted{$n\geq1$},
$M$ est \deleted{\unsure{formé de}} engendré par $\mathbb Q(-n)$\} ; on trouve
en $=$ \unsure{temps} que $n$ est \uline{pair}. Donc on trouve

\diagram{\begin{tikzcd}[ampersand replacement=\&,column sep=small,row sep=small]
(f_\infty=1) \arrow[r,Leftrightarrow] \& (\mathbb R(\alpha)=\mathbb R) \arrow[r,Leftrightarrow] \&
\substack{M\ \text{est engendré par }\mathbb Q(-2n),\\\text{\unsure{pour|un} quelque }n\geq0.} \arrow[d,Leftrightarrow]\\
\& \& G\simeq G_m\ \uline{\text{ou}}\ G\simeq\{e\}
\end{tikzcd}}

% slanted
\margin{\uline{moyennant} \uline{conj. de Hodge}}

\uline{NB} $f_\infty$ ne peut être \uline{central} que si $G^0=\mathrm{Im}\,\deleted{\unsure{G|\varepsilon}}i$, ce qui
implique que $M$ est \deleted{engendré} contenu dans la
$\otimes$-cat. engendrée par $\mathbb Q$ tordu \uline{et} par $\mathbb Q(2)$.
Il revient au $\unsure{=|{-}}$ dire que $G$ est \uline{comm.}.
% archive sheet number: 35
